\section{Refinement Reflection: \corelan}
\label{sec:formalism}
\label{sec:types-reflection}
\label{sec:theory}

Refinement reflection encodes recursive functions
in the quantifier-free, SMT logic and it is
formalized in three steps.
%
First, we develop a core calculus \corelan with an
\emph{undecidable} type system based on denotational
semantics and prove it sound. 
%%We show how the soundness of the type system
%%allows us to \emph{prove theorems} using \corelan.
%
Next, in \S~\ref{sec:theory:algorithmic} we define
a language \smtlan that soundly approximates \corelan
while enabling decidable, SMT-based type checking.
%
Finally, in \S~\ref{sec:pbe} we develop a complete
proof search algorithm to automate equational reasoning.

\subsection{Syntax}

\begin{figure}[t]
\begin{minipage}[t]{0.49\textwidth}
$$
\arraycolsep=2.5pt
\begin{array}{rccl}
\emphbf{Ops}& \odot
  & ::= & = \spmid  <
\\[0.03in]

\emphbf{Consts}& c
  & ::=
  & \land \spmid \lnot \spmid \odot \spmid +,-,\dots  \\
  && \spmid & \etrue\spmid \efalse \spmid 0, -1, 1, \dots
\\[0.03in]

\emphbf{Vals}& w
  & ::=&  c \spmid \efun{x}{\typ}{e} \spmid D\ \overline{w}
\\[0.03in]

\emphbf{Exprs}& e
  & ::=    & w \spmid x \spmid \eapp{e}{e}      \\
  && \spmid & \ecase{x}{e}{\dc}{\overline{x}}{e}
  %% &   & \spmid & \eletb{x}{\gtyp}{e}{e}
\\[0.06in]

\emphbf{Binds}& \bd
  & ::= & e \spmid \eletb{x}{\gtyp}{\bd}{\bd}
\\[0.03in]

\emphbf{Progs}& \prog
  & ::= & \bd \spmid \erefb{x}{\gtyp}{e}{\prog}
\\[0.03in]

\emphbf{Bas. Types}& \btyp
  & ::= & \tint \spmid \tbool \spmid T
\\[0.03in]

\emphbf{Ref. Types}& \typ
  & ::=   &  {\tref{v}{\btyp^{[\tlabel]}}{\refa}} \spmid \trfun{x}{\typ_x}{\typ}
\end{array}
$$

\end{minipage}
\ \
\begin{minipage}[t]{0.49\textwidth}
$$
\arraycolsep=2.5pt
\begin{array}{rccl}
\emphbf{Terms}& \pred
  & ::=    & \pred \binop \pred \spmid \unop \pred \spmid \dc \\
  && \spmid & n \spmid b \spmid x \spmid  x\ \overline{\pred}\\
%%  && \spmid & \forall \overline{\tbind{x}{\sort}}. \pred
  && \spmid & \eif{\pred}{\pred}{\pred}
\\[0.06in]

\emphbf{Ints} & n
  & ::= & 0, -1, 1, \dots
\\[0.03in]

\emphbf{Bools}& b
  & ::= & \etrue \spmid \efalse
\\[0.03in]

\emphbf{Bin-Ops}& \binop
  & ::= & = \spmid < \spmid \land \spmid +, - , \dots
\\[0.03in]

\emphbf{Un-Ops}& \unop
  & ::= & \lnot \spmid \dots
\\[0.03in]

\emphbf{Args}& \sort_a
  & ::= & \tint \spmid \tbool \spmid \tuniv \\
  && \spmid & \tsmtfun{\sort_a}{\sort_a}
\\[0.06in]

\emphbf{Sorts} & \sort
  & ::=    & \sort_a \spmid \sort_a \rightarrow \sort
\end{array}
$$

\end{minipage}

\label{fig:lang-syntax}
\caption{\textbf{(Left) Syntax of \corelan}: Denotational Typing.
         \quad
         \textbf{(Right) Syntax of \smtlan:} Algorithmic Typing.
         }
\label{fig:syntax}
\label{fig:smtsyntax}
\figrule
\end{figure}

%
Figure~\ref{fig:syntax} summarizes the syntax of \corelan,
which is essentially the calculus \undeclang~\cite{Vazou14}
with explicit recursion and a special $\erefname$ binding
to denote terms that are reflected into the refinement logic.
%
The elements of \corelan are constants, values,
expressions, binders, and programs.

\mypara{Constants}
The constants of \corelan
include primitive relations $\odot$,
here, the set $\{ =, <\}$.
%
Moreover, they include the
booleans $\etrue$, $\efalse$,
integers $\mathtt{-1}, \mathtt{0}$, $\mathtt{1}$, \etc,
and logical operators $\mathtt{\land}$, $\mathtt{\lnot}$, \etc.

\mypara{Data Constructors}
%
Data constructors are special constants.
% Each data type has an equality predicate $\haseq{T}$
% that is true only if values of type $T$ can be finitely compared.
For example, the data type \tintlist, which represents
finite lists of integers, has two data constructors: $\dnull$ (nil)
and $\dcons$ (cons).

\mypara{Values \& Expressions}
%
The values of \corelan include
constants, $\lambda$-abstractions
$\efun{x}{\typ}{e}$, and fully
applied data constructors $D$
that wrap values.
%
The expressions of \corelan
include values, variables $x$,
applications $\eapp{e}{e}$, and
$\mathtt{case}$ expressions.

\mypara{Binders \& Programs}
%
A \emph{binder} $\bd$ is a series of possibly recursive
$\mathtt{let}$ definitions, followed by an expression.
%
A \emph{program} \prog is a series of $\erefname$
definitions, each of which names a function
that is reflected into the refinement
logic, followed by a binder.
%
The stratification of programs via binders
is required so that arbitrary recursive definitions
are allowed in the program but cannot be inserted into the logic
via refinements or reflection.
%
(We \emph{can} allow non-recursive $\mathtt{let}$
binders in expressions $e$, but omit them for simplicity.)

\subsection{Operational Semantics}
We define $\hookrightarrow$ to be the small step, call-by-name
$\beta$-reduction semantics for \corelan.
%
We evaluate reflected terms %$\erefb{x}{\gtyp}{e}{\prog}$
as recursive $\mathtt{let}$ bindings, with termination
constraints imposed by the type system:
%
$$
\erefb{x}{\gtyp}{e}{\prog}
\hookrightarrow
\eletb{x}{\gtyp}{e}{\prog}
$$
%
We define $\evalsto{}{}$ to be the reflexive,
transitive closure of $\evals{}{}$.
%
Moreover, we define $\betaeq{}{}$ to be the reflexive,
symmetric, and transitive closure of $\evals{}{}$.

\mypara{Constants} Application of a constant requires the
argument be reduced to a value; in a single step, the
expression is reduced to the output of the primitive
constant operation, \ie ${c\ v} \hookrightarrow \ceval{c}{v}$.
%
For example, consider $=$, the primitive equality
operator on integers.
%
We have $\ceval{=}{n} \defeq =_n$
where $\ceval{=_n}{m}$ equals \etrue
iff $m$ is the same as $n$.
%
%%\VC{This is the first appearance of
%%  $\ceval{c}{v}$, I don't know what it is,
%%  it looks like a relation. Same for
%%  $=_n$.}
%%\NV{This is definition by example, since delta is standard in operational semantics}
%

\mypara{Equality}
We assume that the equality operator
is defined \emph{for all} values,
and, for functions, is defined as
extensional equality.
%
%%\VC{Do we need to assume it, or we could
%%  just define it, since we know what all
%%  those values are?}
%%\NV{You cannot really define operational function equality}
%
That is, for all
$f$ and $f'$,
$\evals{(f = f')}{\etrue}$
   $\mbox{iff}$
  $\forall v.\ \betaeq{f\ v}{f'\ v}$.
%
We assume source \emph{terms} only contain implementable equalities
over non-function types; while function extensional equality only
appears in \emph{refinements}. % and is approximated by the underlying logic.
\VC{This statement is very unclear to me.}

\subsection{Types}

\corelan types include basic types, which are \emph{refined} with predicates,
and dependent function types.
%
\emph{Basic types} \btyp comprise integers, booleans, and a family of data-types
$T$ (representing lists, trees \etc).
%
For example, the data type \tintlist represents lists of integers.
%
We refine basic types with predicates (boolean-valued expressions \refa) to obtain
\emph{basic refinement types} $\tref{v}{\btyp}{\refa}$.
%
%
We use \tlabel to mark provably terminating
computations and use
refinements to ensure that if
${{e}\text{:}{\tref{v}{\btyp^\tlabel}{\refa'}}}$,
then $e$ terminates.
%
As discussed by~\citet{Vazou14} termination labels
are checked using refinement types and are used
to ensure that refinements cannot diverge as required
for soundness of type checking under lazy evaluation.
%
Termination checking is crucial for this work,
as combined with checks for exhaustive definitions (\S~\ref{subsection:overview:refinements}),
it ensures totality (well-formedness)
of expressions as required by
propositions-as-types (\S~\ref{sec:higher-order})
and termination of \pbesym (\S~\ref{sec:pbe}).
%
Finally, we have dependent \emph{function types} $\trfun{x}{\typ_x}{\typ}$
where the input $x$ has the type $\typ_x$ and the output $\typ$ may
refer to the input binder $x$.
%
We write $\btyp$ to abbreviate $\tref{v}{\btyp}{\etrue}$
and \tfunbasic{\typ_x}{\typ} to abbreviate \trfun{x}{\typ_x}{\typ}, if
$x$ does not appear in $\typ$.
%
\VC{Kind of hard to read, even though $x$
  isn't bound, the type still has $x$ as a
  subscript. Compare to, $a:A \to B$, and
  $A \to B$}

\mypara{Denotations}
%
Each type $\typ$ \emph{denotes} a set of expressions $\interp{\typ}$,
that is defined via the operational semantics~\cite{Knowles10}.
%
Let $\shape{\typ}$ be the type we get if we erase all refinements
from $\typ$ and $\bhastype{}{e}{\shape{\typ}}$ be the
standard typing relation for the typed lambda calculus.
%
Then, we define the denotation of types as:
%
\begin{align*}
\interp{\tref{x}{\btyp}{r}} \defeq &\
  \{ e \mid  \bhastype{}{e}{\btyp}
           , \mbox{ if } \evalsto{e}{w} \mbox{ then } \evalsto{r\subst{x}{w}}{\etrue}
  \}\\
\interp{\tref{x}{\btyp^\tlabel}{r}} \defeq &\
  \interp{\tref{x}{\btyp}{r}} \cap \{ e \mid  \exists\ w. \evalsto{e}{w}\}\\
\interp{\trfun{x}{\typ_x}{\typ}} \defeq &\
  \{ e \mid \bhastype{}{e}{\shape{\tfunbasic{\typ_x}{\typ}}}
           , \forall e_x \in \interp{\typ_x}.\ (\eapp{e}{e_x}) \in \interp{\typ\subst{x}{e_x}}
  \}
\end{align*}
%

\mypara{Constants}
For each constant $c$ we define its type \constty{c}
such that $c \in \interp{\constty{c}}$.
%
For example,
%
$$
\begin{array}{lcl}
\constty{3} &\doteq& \tref{v}{\tint^\tlabel}{v = 3}\\
\constty{+} &\doteq& \trfun{\ttx}{\tint^\tlabel}{\trfun{\tty}{\tint^\tlabel}{\tref{v}{\tint^\tlabel}{v = x + y}}}\\
\constty{\leq} &\doteq& \trfun{\ttx}{\tint^\tlabel}{\trfun{\tty}{\tint^\tlabel}{\tref{v}{\tbool^\tlabel}{v \Leftrightarrow x \leq y}}}\\
\end{array}
$$
\VC{The right hand side is now a
  refinement type, not a set like the
  previous one.}

\subsection{Refinement Reflection}\label{subsec:logicalannotations}
%
Reflection \emph{strengthens} function output types
with a refinement that \emph{reflects} the
definition of the function in the logic.
%
We do this by treating each
%
$\erefname$-binder
%
(${\erefb{f}{\gtyp}{e}{\prog}}$)
%
as a $\eletname$-binder
%
(${\eletb{f}{\exacttype{\gtyp}{e}}{e}{\prog}}$)
%
during type checking (rule $\rtreflect$ in Figure~\ref{fig:typing}).

\mypara{Reflection}
%
We write \exacttype{\typ}{e} for the \emph{reflection}
of the term $e$ into the type $\typ$, defined as % by strengthening \typ as:
%
$$
\begin{array}{lcl}
\exacttype{\tref{v}{\btyp{^\tlabel}}{r}}{e}
  & \defeq
  & \tref{v}{\btyp{^\tlabel}}{r \land v = e}\\
\exacttype{\trfun{x}{\typ_x}{\typ}}{\efun{x}{}{e}}
  & \defeq
  & \trfun{x}{\typ_x}{\exacttype{\typ}{e}}
\end{array}
$$
%
As an example, recall from \S~\ref{sec:overview}
that the "reflect fib" strengthens the type of
"fib" with the refinement "fibP".
%
That is, let the user specified type of \fibname
be $t_\fibnameblack$ and the its definition be definition
$\lambda n . e_\fibnameblack$.
%
\begin{align*}
t_\fibnameblack & \defeq {\tfun{\ttnat}{\ttnat}} \\
e_\fibnameblack & \defeq \efib
\end{align*}
%
Then, the reflected type of \fibname will be:
%
$$
\exacttype{t_\fibnameblack}{e_\fibnameblack} = \trfun{n}{\ttnat}{\tref{v}{\tint^\tlabel}{0 \leq v \land v = e_\fibname}}
$$


\mypara{Termination Checking}
We defined \exacttype{\cdot}{\cdot} to be a \textit{partial}
function that only reflects provably terminating expressions,
\ie expressions whose result type is marked with \tlabel.
%
If a non-provably terminating function is reflected
in an \corelan expression then type checking will
fail (with a reflection type error in the implementation).
%
This restriction is crucial for soundness,
as diverging expressions can lead to inconsistencies.
%
For example, reflecting the diverging "f x = 1 + f x"
into the logic leads to an inconsistent system that
is able to prove "0 = 1".
%%This choice complies with both the facts that
%%all refinement expressions are checked to be terminating
%%by the type system (by the rule \rwbase) and, as explained in \S~\ref{sec:algorithmic},
%%all reflected functions are encoded as total uninterpreted functions
%%in the SMT logic of \smtlan.

\mypara{Automatic Reflection}
Reflection of \corelan expressions into the refinements happens
automatically by the type system, not manually by the user.
%
The user simply annotates a function $f$ as $\annotReflect\ f$.
%
Then, the rule $\rtreflect$ in Figure~\ref{fig:typing}
is used to type check the reflected function by strengthening
$f$'s result via \exacttype{\cdot}{\cdot}.
%
Finally, the rule \rtlet is used to check that the
automatically strengthened type of $f$ satisfies $f$'s implementation.

\RN{Hmm, the ordering is a bit weird here where there are a couple references to
fig 5, but then the next subsec starts with an introductory tone.}

\subsection{Typing Rules}
\begin{figure}[!t]
%% \begin{minipage}[t]{0.47\textwidth}
\emphbf{Typing}\hfill{\fbox{\hastype{\env}{\RRenv}{\prog}{\typ}}}\\
$$
\inference{
	\tbind{x}{\gtyp}\in\env
}{
	\hastype{\env}{\RRenv}{x}{\gtyp}
}[\rtvar]
\qquad
\inference{
}{
	\hastype{\env}{\RRenv}{c}{\constty{c}}
}[\rtconst]
\quad
\inference{
	\hastype{\env}{\RRenv}{\prog}{\typ'} &
	\issubtype{\env}{\RRenv}{\typ'}{\typ}
}{
	\hastype{\env}{\RRenv}{\prog}{\typ}
}[\rtsub]
$$

$$
\inference{
	\hastype{\env}{\RRenv}{e}{\tref{v}{\btyp}{e_r}}
}{
	\hastype{\env}{\RRenv}{e}{\tref{v}{\btyp}{e_r\land v = e}}
}[\rtexact]
\qquad
\inference{
	\hastype{\env, \tbind{x}{\typ_x}}{\RRenv}{e}{\typ}
}{
	\hastype{\env}{\RRenv}{\efun{x}{\typ}{e}}{(\trfun{x}{\typ_x}{\typ})}
}[\rtfun]
$$

$$
\inference{
	\hastype{\env}{\RRenv}{e}{(\trfun{x}{\typ_x}{\typ})} &&
	\hastype{\env}{\RRenv}{e_x}{\typ_x}
}{
	\hastype{\env}{\RRenv}{e\ e_x}{\typ\subst{x}{e_x}}
}[\rtapp]
\qquad
\inference
	{\hastype{\env, \tbind{x}{\gtyp_x}}{\RRenv}{\bd_x}{\gtyp_x} &
	 \iswellformed{\env, \tbind{x}{\gtyp_x}}{\typ_x} \\
	 \hastype{\env, \tbind{x}{\gtyp_x}}{\RRenv}{\bd}{\gtyp} &
	 \iswellformed{\env}{\typ}}
	{\hastype{\env}{\RRenv}{\eletb{x}{\gtyp_x}{\bd_x}{\bd}}{\typ}}
	[\rtlet]
$$

$$
\inference{
	\hastype{\env}{\RRenv}{e}{\tref{v}{T}{e_r}} & \iswellformed{\env}{\typ} \\
	& \forall i. \constty{\dc_i} = \tfunbasic{\overline{\tbind{y_j}{\typ_j}}}{\tref{v}{T}{e_{r_i}}} 
	& \hastype{\env, \overline{\tbind{y_j}{\typ_j}}, \tbind{x}{\tref{v}{T}{e_r \land e_{r_i}}} }{\RRenv}{e_i}{\typ}
}{
	\hastype{\env}{\RRenv}{\ecase{x}{e}{\dc_i}{\overline{y_i}}{e_i}}{\typ}
}[\rtcase]
$$

$$
\inference
	{\hastype{\env}{\RRenv, f \mapsto e}
	 				 {\eletb{f}{\exacttype{\gtyp_f}{e}}{e}{\prog}}
					 {\typ}
	}
	{\hastype{\env}{\RRenv}
					 {\erefb{f}{\gtyp_f}{e}{\prog}}
					 {\typ}
	}[\rtreflect]
$$

%% \end{minipage}
%% \hspace{0.2in}
%% \begin{minipage}[t]{0.47\textwidth}
\emphbf{Well Formedness}\hfill{\fbox{\iswellformed{\env}{\typ}}}\\

$$
\inference{
  \hastype{\env,\tbind{v}{\btyp}}{\emptyset}{\refa}{\tbool^{\tlabel}}
}{
	\iswellformed{\env}{\tref{v}{\btyp}{\refa}}
}[\rwbase]
\qquad
\inference{
	\iswellformed{\env}{\typ_x} &
	\iswellformed{\env,\tbind{x}{\typ_x}}{\typ}
}{
	\iswellformed{\env}{\trfun{x}{\typ_x}{\typ}}
}[\rwfun]
$$

\emphbf{Subtyping}\hfill{\fbox{\issubtype{\env}{\RRenv}{\typ_1}{\typ_2}}}\\

$$
\inference{
        \forall \theta\in\interp{\env}.
        \interp{\applysub{\theta}{\tref{v}{\btyp}{\refa_1}}}
        \subseteq
        \interp{\applysub{\theta}{\tref{v}{\btyp}{\refa_2}}}
}{
        \issubtype{\env}{\RRenv}{\tref{v}{\btyp}{\refa_1}}{\tref{v}{\btyp}{\refa_2}}
}[\rsubbasecore]
$$

$$
\inference{
	\issubtype{\env}{\RRenv}{\typ_x'}{\typ_x} &
	\issubtype{\env,\tbind{x}{\typ_x'}}{\RRenv}{\typ}{\typ'}
}{
	\issubtype{\env}{\RRenv}{\trfun{x}{\typ_x}{\typ}}{\trfun{x}{\typ_x'}{\typ'}}
}[\rsubfun]
$$
%%\end{minipage}
\caption{Typing of \corelan.} % and algorithmic subtyping of \smtlan}
\label{fig:typing}
\figrule
\end{figure}

%
Next, we present the type-checking rules of \corelan, as found in Figure~\ref{fig:typing}.

\mypara{Environments and Closing Substitutions}
A \emph{type environment} $\env$ is a sequence of type bindings
$\tbind{x_1}{\typ_1},\ldots,\tbind{x_n}{\typ_n}$. An environment
denotes a set of \emph{closing substitutions} $\sto$ which are
sequences of expression bindings:
$\gbind{x_1}{e_1}, \ldots,$ $\gbind{x_n}{e_n}$ such that:
$$
\interp{\env} \defeq  \{\sto \mid \forall \tbind{x}{\typ} \in \Env.
                                    \sto(x) \in \interp{\applysub{\sto}{\typ}} \}
$$
where $\applysub{\sto}{\typ}$ applies a substitution to a type (and
  likewise $\applysub{\sto}{\prog}$, to a program).

A reflection environment $\RRenv$ is a sequence that binds the names
of the reflected functions with their definitions
$\gbind{f_1}{e_1},\ldots,\gbind{f_n}{e_n}$.
%
A reflection environment respects a type environment
when all reflected functions satisfy their types:
$$
\env \models \RRenv \defeq
\forall (\gbind{f}{e}) \in \RRenv.\;
\exists \typ.\;
(\tbind{f}{\typ}) \in \env \wedge
(\hastype{\env}{\RRenv}{e}{\typ})
$$

\mypara{Typing}
A judgment \hastype{\env}{\RRenv}{\prog}{\typ} states that
the program $\prog$ has the type $\typ$ in
the type environment $\env$ and the reflection environment $\RRenv$.
That is, when the free variables in $\prog$ are
bound to expressions described by $\env$, the
program $\prog$ will evaluate to a value
described by $\typ$.

\mypara{Rules}
%
All but two of the rules are standard~\cite{Knowles10,Vazou14}
except for the addition of the reflection environment $\RRenv$ at each rule.
%
First, rule \rtreflect is used to
extend the reflection environment with the binding of the function name
with its definition ($f \mapsto e$)
and moreover to strengthen the type of each
reflected binder with its definition, as described previously
in \S~\ref{subsec:logicalannotations}.
%
Second, rule \rtexact strengthens the expression with
a singleton type equating the value and the expression
(\ie reflecting the expression in the type).
%
This is a generalization of the ``selfification'' rules
from \cite{Ou2004,Knowles10} and is required to
equate the reflected functions with their definitions.
%
For example, the application "fib 1" is typed as
${\tref{v}{\tint^\tlabel}{\fibdef\ 1 \wedge v = \fib\ 1}}$ where
the first conjunct comes from the (reflection-strengthened)
output refinement of "fib" (\S~\ref{sec:overview}) and
the second comes from rule~\rtexact.

\mypara{Well-formedness}
%
A judgment \iswellformed{\env}{\typ} states that
the refinement type $\typ$ is well-formed in
the environment $\env$.
%
Following~\citet{Vazou14}, $\typ$ is well-formed if all
the refinements in $\typ$ are $\tbool$-typed,
provably terminating expressions in $\env$.

%%\mypara{Termination}
%%%
%%Under arbitrary beta-reduction semantics
%%(which includes lazy evaluation), soundness
%%of refinement type checking requires checking
%%termination, for two reasons:
%%%
%%(1)~to ensure that refinements cannot diverge, and
%%(2)~to account for the environment during subtyping~\citep{Vazou14}.
%%%
%%We use \tlabel to mark provably terminating
%%computations, and extend the rules to use
%%refinements to ensure that if
%%${\ahastype{\env}{e}{\tref{v}{\btyp^\tlabel}{r}}}$,
%%then $e$ terminates~\citep{Vazou14}.

\mypara{Subtyping}
%
A judgment \issubtype{\env}{\RRenv}{\typ_1}{\typ_2} states
that the type $\typ_1$ is a subtype of %the type
$\typ_2$ in the environments $\env$ and $\RRenv$.
%
%%Informally, $\typ_1$ is a subtype of $\typ_2$ if
%%the refinement of $\typ_1$ \emph{implies}
%%the refinement of $\typ_2$
%%under the assumptions described by $\env$.
%
Informally, $\typ_1$ is a subtype of $\typ_2$ if, when
the free variables of $\typ_1$ and $\typ_2$
are bound to expressions described by $\env$,
the denotation of $\typ_1$
is \emph{contained in} the denotation of $\typ_2$.
%
Subtyping of basic types reduces to denotational
containment checking, shown in rule \rsubbasecore.
%
That is, $\typ_1$ is a subtype of $\typ_2$ under $\env$
if for any closing substitution $\sto$ in $\interp{\env}$,
$\interp{\applysub{\sto}{\typ_1}}$ is contained in
$\interp{\applysub{\sto}{\typ_2}}$.

\mypara{Soundness}
%
We prove that 
typing implies denotational
inclusion and 
evaluation preserves typing.
%
\begin{theorem}{[Soundness of \corelan]}\label{thm:safety}
\begin{itemize}
\item\textbf{Denotations}
If $\;\hastype{\env}{\RRenv}{\prog}{\typ}$, then
$\forall \sto\in \interp{\env}. \applysub{\sto}{\prog} \in \interp{\applysub{\sto}{\typ}}$.
\item\textbf{Preservation}
If \;\hastype{\emptyset}{\emptyset}{\prog}{\typ}
       and $\evalsto{\prog}{w}$, then $\hastype{\emptyset}{\emptyset}{w}{\typ}$.
\end{itemize}
\end{theorem}

The proofs can be found in~\cite{appendix}.
Theorem~\ref{thm:safety} lets us
prove that if \formula is a \corelan type
interpreted as a proposition
(using the mapping of Figure~\ref{fig:hol})
and if there exists a $\prog$ so that
$\hastype{\emptyset}{\emptyset}{\prog}{\formula}$,
the \formula is valid.
%%Theorem~\ref{thm:safety} lets us interpret well
%%typed programs as proofs of propositions.
%
%%\VC{Does the subsumption rule make \corelan\ undecidable?}
%%\NV{Yes}
%
For example, in \S~\ref{sec:overview} we verified
that the term \fibincrname proves
${\trfun{n}{\tnat}{\ttref{\fib{n} \leq \fib{(n+1)}}}}$.
%
Via soundness of \corelan, we get that for each valid input "n",
the result refinement is valid.
$$
\forall n. \evalsto{0 \leq n}{\etrue} \Rightarrow \evalsto{\texttt{fib}\ n \leq \texttt{fib}\ {(n+1)}}{\etrue}
$$
%
